% =============================================================================
% Chapter 8: The Complete GPT Model
% =============================================================================
\chapter{The Complete GPT Model}
\label{ch:complete-gpt}

\begin{objectives}
    \item Understand the complete GPT architecture
    \item Trace data flow from input to output
    \item Learn how loss is computed for language modeling
    \item Master weight initialization strategies
    \item Calculate exact parameter counts
\end{objectives}

% -----------------------------------------------------------------------------
\section{GPT Architecture Overview}
% -----------------------------------------------------------------------------

We've now covered all the components. Let's assemble them into the complete GPT model.

\begin{figure}[H]
\centering
\begin{tikzpicture}[
    box/.style={rectangle, draw, rounded corners, minimum width=4cm, minimum height=0.8cm, align=center, fill=blue!10},
    block/.style={rectangle, draw, rounded corners, minimum width=4cm, minimum height=1.5cm, align=center, fill=orange!10},
    arrow/.style={-{Stealth[length=2mm]}, thick},
]
    % Input
    \node[box, fill=green!10] (input) at (0, 0) {Token IDs\\$[B, S]$};

    % Embedding
    \node[box, fill=purple!20] (emb) at (0, -1.5) {Token Embedding\\$[B, S] \to [B, S, 384]$};

    % Dropout
    \node[box, fill=gray!20] (drop) at (0, -2.8) {Dropout};

    % Transformer blocks
    \node[block] (blocks) at (0, -4.5) {Transformer Blocks $\times 6$\\Attention + FFN + Residuals};

    % Final norm
    \node[box, fill=yellow!20] (norm) at (0, -6.3) {RMSNorm (final)};

    % LM Head
    \node[box, fill=purple!20] (head) at (0, -7.6) {LM Head (weight-tied)\\$[B, S, 384] \to [B, S, 16384]$};

    % Output
    \node[box, fill=red!10] (output) at (0, -9.1) {Logits\\$[B, S, V]$};

    % Arrows
    \draw[arrow] (input) -- (emb);
    \draw[arrow] (emb) -- (drop);
    \draw[arrow] (drop) -- (blocks);
    \draw[arrow] (blocks) -- (norm);
    \draw[arrow] (norm) -- (head);
    \draw[arrow] (head) -- (output);

    % Dimension annotations
    \node[right=0.5cm of input, font=\footnotesize, gray] {batch $\times$ seq\_len};
    \node[right=0.5cm of emb, font=\footnotesize, gray] {+ d\_model dim};
    \node[right=0.5cm of head, font=\footnotesize, gray] {project to vocab};
\end{tikzpicture}
\caption{Complete GPT architecture: from token IDs to logits}
\label{fig:gpt-architecture}
\end{figure}

% -----------------------------------------------------------------------------
\section{The GPT Class}
% -----------------------------------------------------------------------------

Here's the complete model implementation:

\begin{lstlisting}[style=python, caption={GPT model from model.py}]
class GPT(nn.Module):
    def __init__(self, config: ModelConfig) -> None:
        super().__init__()
        self.config = config

        # Token embeddings
        self.token_embedding = nn.Embedding(
            config.vocab_size,    # 16384
            config.d_model,       # 384
        )

        # Dropout
        self.dropout = nn.Dropout(config.dropout)

        # Stack of transformer blocks
        self.blocks = nn.ModuleList([
            TransformerBlock(config)
            for _ in range(config.n_layers)  # 6
        ])

        # Final layer norm
        self.final_norm = RMSNorm(config.d_model)

        # Output projection (language model head)
        self.lm_head = nn.Linear(
            config.d_model,       # 384
            config.vocab_size,    # 16384
            bias=False,
        )

        # Weight tying
        self.lm_head.weight = self.token_embedding.weight

        # Initialize weights
        self.apply(self._init_weights)
\end{lstlisting}

% -----------------------------------------------------------------------------
\section{The Forward Pass}
% -----------------------------------------------------------------------------

Let's trace exactly what happens when data flows through the model.

\subsection{Input Processing}

\begin{lstlisting}[style=python, caption={Forward pass from model.py}]
def forward(
    self,
    input_ids: Tensor,              # [batch, seq_len]
    attention_mask: Tensor | None = None,
    labels: Tensor | None = None,
) -> dict[str, Tensor]:

    batch_size, seq_len = input_ids.shape

    # Step 1: Token embedding
    x = self.token_embedding(input_ids)  # [B, S, 384]
    x = self.dropout(x)
\end{lstlisting}

\subsection{Transformer Blocks}

\begin{lstlisting}[style=python]
    # Step 2: Pass through all transformer blocks
    for block in self.blocks:
        x = block(x, attention_mask)
    # Still [B, S, 384]
\end{lstlisting}

\subsection{Output Projection}

\begin{lstlisting}[style=python]
    # Step 3: Final normalization
    x = self.final_norm(x)  # [B, S, 384]

    # Step 4: Project to vocabulary
    logits = self.lm_head(x)  # [B, S, 16384]
\end{lstlisting}

\subsection{Loss Computation}

\begin{lstlisting}[style=python]
    # Step 5: Compute loss if labels provided
    if labels is not None:
        # Shift logits and labels for next-token prediction
        shift_logits = logits[..., :-1, :].contiguous()
        shift_labels = labels[..., 1:].contiguous()

        # Cross-entropy loss
        loss = F.cross_entropy(
            shift_logits.view(-1, self.config.vocab_size),
            shift_labels.view(-1),
            ignore_index=-100,  # Ignore padding
        )
        return {"logits": logits, "loss": loss}

    return {"logits": logits}
\end{lstlisting}

% -----------------------------------------------------------------------------
\section{Understanding the Loss Computation}
% -----------------------------------------------------------------------------

The loss computation involves a shift that aligns predictions with targets.

\subsection{The Shift Operation}

\begin{figure}[H]
\centering
\begin{tikzpicture}[
    box/.style={rectangle, draw, minimum width=1.2cm, minimum height=0.6cm, font=\small},
]
    % Logits
    \node at (-2, 1) {\textbf{Logits:}};
    \foreach \x/\t in {0/0, 1.3/1, 2.6/2, 3.9/3, 5.2/4} {
        \node[box, fill=blue!10] at (\x, 1) {$\ell_{\t}$};
    }

    % Shifted logits (for loss)
    \node at (-2, 0) {\textbf{Shift logits:}};
    \foreach \x/\t in {0/0, 1.3/1, 2.6/2, 3.9/3} {
        \node[box, fill=green!20] at (\x, 0) {$\ell_{\t}$};
    }
    \node[box, fill=gray!30] at (5.2, 0) {---};

    % Labels
    \node at (-2, -1) {\textbf{Labels:}};
    \foreach \x/\t in {0/0, 1.3/1, 2.6/2, 3.9/3, 5.2/4} {
        \node[box, fill=orange!10] at (\x, -1) {$y_{\t}$};
    }

    % Shifted labels (for loss)
    \node at (-2, -2) {\textbf{Shift labels:}};
    \node[box, fill=gray!30] at (0, -2) {---};
    \foreach \x/\t in {1.3/1, 2.6/2, 3.9/3, 5.2/4} {
        \node[box, fill=orange!20] at (\x, -2) {$y_{\t}$};
    }

    % Alignment arrows
    \foreach \x in {0, 1.3, 2.6, 3.9} {
        \draw[-{Stealth}, gray] (\x, -0.35) -- (\x + 1.3, -1.65);
    }

    % Annotation
    \node[font=\footnotesize, gray] at (2.6, -3) {Logit at position $i$ predicts label at position $i+1$};
\end{tikzpicture}
\caption{Shift alignment: each logit predicts the next token}
\label{fig:loss-shift}
\end{figure}

\subsection{Why Shift?}

Remember the language modeling objective: predict the \textit{next} token. So:
\begin{itemize}
    \item Logits at position 0 should predict token at position 1
    \item Logits at position $n-2$ should predict token at position $n-1$
    \item Logits at position $n-1$ have nothing to predict (discarded)
\end{itemize}

\subsection{Cross-Entropy Loss}

The cross-entropy loss measures how well predictions match targets:

\begin{equation}
\mathcal{L} = -\frac{1}{N} \sum_{i=1}^{N} \log p(y_i | x_{<i})
\end{equation}

where $p(y_i | x_{<i})$ is the predicted probability of the correct next token.

\begin{keypoint}
\textbf{Perplexity} is $\exp(\mathcal{L})$---the exponential of the average loss. Lower perplexity means better predictions. A perplexity of 10 means the model is, on average, as ``surprised'' as if choosing from 10 equally likely options.
\end{keypoint}

% -----------------------------------------------------------------------------
\section{Weight Initialization}
% -----------------------------------------------------------------------------

Proper initialization is crucial for stable training.

\subsection{TryLM's Initialization Strategy}

\begin{lstlisting}[style=python, caption={Weight initialization from model.py}]
def _init_weights(self, module: nn.Module) -> None:
    if isinstance(module, nn.Linear):
        torch.nn.init.normal_(module.weight, mean=0.0, std=0.02)
        if module.bias is not None:
            torch.nn.init.zeros_(module.bias)

    elif isinstance(module, nn.Embedding):
        torch.nn.init.normal_(module.weight, mean=0.0, std=0.02)
\end{lstlisting}

\subsection{Scaled Initialization for Residuals}

Output projections get special treatment:

\begin{lstlisting}[style=python]
# Scale down residual projections
for name, p in self.named_parameters():
    if name.endswith("out_proj.weight") or name.endswith("down_proj.weight"):
        torch.nn.init.normal_(
            p,
            mean=0.0,
            std=0.02 / math.sqrt(2 * config.n_layers)
        )
\end{lstlisting}

\subsection{Why Scale?}

Residual connections sum contributions from many layers. Without scaling, the variance grows with depth:

\begin{equation}
\text{Variance after } L \text{ layers} \propto L
\end{equation}

Scaling by $1/\sqrt{2L}$ keeps variance bounded, enabling stable training of deep networks.

% -----------------------------------------------------------------------------
\section{Parameter Count}
% -----------------------------------------------------------------------------

Let's calculate the exact parameter count for \TryLM.

\subsection{Component Breakdown}

\begin{table}[H]
\centering
\begin{tabular}{lrr}
\toprule
\textbf{Component} & \textbf{Formula} & \textbf{Parameters} \\
\midrule
Token Embedding & $V \times d$ & $16384 \times 384 = 6,291,456$ \\
\addlinespace
\textit{Per Transformer Block:} & & \\
\quad QKV Projection & $d \times 3d$ & $384 \times 1152 = 442,368$ \\
\quad Output Projection & $d \times d$ & $384 \times 384 = 147,456$ \\
\quad Gate Projection & $d \times h$ & $384 \times 1024 = 393,216$ \\
\quad Up Projection & $d \times h$ & $384 \times 1024 = 393,216$ \\
\quad Down Projection & $h \times d$ & $1024 \times 384 = 393,216$ \\
\quad 2 × RMSNorm & $2 \times d$ & $2 \times 384 = 768$ \\
\quad \textit{Block Total} & & $1,770,240$ \\
\addlinespace
6 Blocks & $6 \times 1,770,240$ & $10,621,440$ \\
Final RMSNorm & $d$ & $384$ \\
LM Head & (tied with embedding) & $0$ \\
\midrule
\textbf{Total} & & $\mathbf{16,913,280}$ \\
\bottomrule
\end{tabular}
\caption{Detailed parameter breakdown for TryLM}
\label{tab:parameter-count}
\end{table}

\begin{note}
The actual count may vary slightly due to rounding in the FFN hidden dimension. The model is approximately 10.8M parameters after weight tying (since embeddings are counted once).
\end{note}

\subsection{Parameter Count Method}

\begin{lstlisting}[style=python, caption={Parameter counting from config.py}]
def count_parameters(self) -> dict[str, int]:
    """Estimate parameter counts for each component."""
    d = self.d_model
    h = int(2 * self.d_ff / 3)
    h = 256 * ((h + 255) // 256)  # Round to 256

    token_emb = self.vocab_size * d
    block_params = (
        3 * d * d +      # QKV projection
        d * d +          # Output projection
        3 * d * h +      # FFN: gate, up, down
        2 * d            # 2 RMSNorm layers
    )
    total_blocks = block_params * self.n_layers
    final_norm = d

    return {
        "token_embedding": token_emb,
        "per_block": block_params,
        "all_blocks": total_blocks,
        "final_norm": final_norm,
        "lm_head": 0,  # Tied with embedding
        "total": token_emb + total_blocks + final_norm,
    }
\end{lstlisting}

% -----------------------------------------------------------------------------
\section{Model Configuration}
% -----------------------------------------------------------------------------

The model is fully configured through the \texttt{ModelConfig} dataclass:

\begin{lstlisting}[style=python, caption={ModelConfig from config.py}]
@dataclass
class ModelConfig:
    vocab_size: int = 16384      # Vocabulary size
    n_layers: int = 6            # Number of transformer blocks
    n_heads: int = 6             # Number of attention heads
    d_model: int = 384           # Model/embedding dimension
    d_ff: int = 1536             # FFN hidden dimension
    context_length: int = 512    # Maximum sequence length
    dropout: float = 0.1         # Dropout rate
    bias: bool = False           # Use bias in linear layers

    @property
    def head_dim(self) -> int:
        return self.d_model // self.n_heads  # 384 / 6 = 64
\end{lstlisting}

% -----------------------------------------------------------------------------
\section{Model Variants}
% -----------------------------------------------------------------------------

You can easily create different model sizes by adjusting the configuration:

\begin{table}[H]
\centering
\begin{tabular}{lrrrrr}
\toprule
\textbf{Variant} & \textbf{Layers} & \textbf{Heads} & \textbf{$d$} & \textbf{FFN} & \textbf{Params} \\
\midrule
Tiny & 4 & 4 & 256 & 1024 & $\sim$5M \\
Small (default) & 6 & 6 & 384 & 1536 & $\sim$11M \\
Medium & 8 & 8 & 512 & 2048 & $\sim$25M \\
Large & 12 & 12 & 768 & 3072 & $\sim$85M \\
\bottomrule
\end{tabular}
\caption{Different model sizes by configuration}
\label{tab:model-variants}
\end{table}

\begin{lstlisting}[style=python]
# Create a medium model
config = ModelConfig(
    n_layers=8,
    n_heads=8,
    d_model=512,
    d_ff=2048,
)
model = GPT(config)
\end{lstlisting}

% -----------------------------------------------------------------------------
\section{Putting It All Together}
% -----------------------------------------------------------------------------

Here's a complete example of creating and using the model:

\begin{lstlisting}[style=python]
from trylm.model import GPT
from trylm.config import ModelConfig
from trylm.tokenizer import TryLMTokenizer
import torch

# Load tokenizer
tokenizer = TryLMTokenizer.load("data/tokenizer")

# Create model
config = ModelConfig(vocab_size=tokenizer.vocab_size)
model = GPT(config)
model.eval()

# Prepare input
text = "Once upon a time"
input_ids = torch.tensor([tokenizer.encode(text)])  # [1, seq_len]

# Forward pass (inference)
with torch.no_grad():
    outputs = model(input_ids)
    logits = outputs["logits"]  # [1, seq_len, vocab_size]

# Get prediction for next token
next_token_logits = logits[0, -1, :]  # [vocab_size]
next_token_id = next_token_logits.argmax().item()
next_token = tokenizer.decode([next_token_id])
print(f"Predicted next token: {next_token}")
\end{lstlisting}

% -----------------------------------------------------------------------------
\section{Exercises}
% -----------------------------------------------------------------------------

\begin{exercise}
\textbf{Trace Tensor Shapes}

For batch\_size=4 and seq\_len=128, trace the shape through each step:
\begin{enumerate}
    \item After token embedding
    \item After each transformer block
    \item After final norm
    \item After lm\_head
\end{enumerate}
\end{exercise}

\begin{exercise}
\textbf{Verify Weight Tying}

Check that weight tying works correctly:
\begin{lstlisting}[style=python, numbers=none]
model = GPT(ModelConfig())

# These should be the same tensor
print(model.token_embedding.weight.data_ptr())
print(model.lm_head.weight.data_ptr())

# Modify one, check the other changes
model.token_embedding.weight[0, 0] = 999.0
print(model.lm_head.weight[0, 0])  # Should be 999.0
\end{lstlisting}
\end{exercise}

\begin{exercise}
\textbf{Calculate Memory Usage}

For a batch\_size=32 and seq\_len=512 forward pass:
\begin{enumerate}
    \item Memory for model parameters (float32)
    \item Memory for attention scores (all 6 heads, all layers)
    \item Memory for activations (estimate)
\end{enumerate}
\end{exercise}

\begin{exercise}
\textbf{Create a Custom Model}

Define a custom model configuration with:
\begin{itemize}
    \item 4 layers
    \item 4 heads
    \item 256 model dimension
    \item Context length of 256
\end{itemize}

Calculate its parameter count and verify with \texttt{model.count\_parameters()}.
\end{exercise}

% -----------------------------------------------------------------------------
\section{Summary}
% -----------------------------------------------------------------------------

\begin{summary}
    \item The \textbf{GPT model} = embedding + transformer blocks + norm + output projection
    \item \textbf{Forward pass}: tokens $\to$ embeddings $\to$ blocks $\to$ norm $\to$ logits
    \item \textbf{Loss computation} uses a shift to align predictions with next-token targets
    \item \textbf{Weight initialization} uses $\mathcal{N}(0, 0.02)$ with scaled residual projections
    \item \textbf{Weight tying} shares parameters between embedding and output projection
    \item \TryLM has $\sim$\textbf{10.8M parameters} (after tying)
    \item Model size is easily \textbf{configurable} through \texttt{ModelConfig}
\end{summary}

We now have a complete understanding of the GPT architecture! In Part III, we'll learn how to train this model to actually generate coherent text.
